\chapter{Problem, Fragestellung, Vision}
\label{cha:problem_fragestellung_vision}

\section{Problemstellung}

Das \gls{tor}-Netzwerk anonymisiert Verbindungen, indem Pakete über mehrere verschlüsselte Relays geleitet werden.
Trotz dieses Schutzes zeigen \gls{wf}-Angriffe, dass ein passiver Beobachter zwischen Client Eingangsknoten besuchte Webseiten allein anhand der Paketrichtungen und Zeitverläufe zuordnen kann.
Mit dem Deep-Learning-Verfahren \gls{df} werden in Closed-World-Studien Erkennungsraten von über 95~\% berichtet \cite{sirinam2018deep}.

Als Gegenmassnahme integriert das Tor-Projekt seit Version 0.4 ein Circuit-Padding-Framework, das zusätzliche Zellen in den Verkehrsstrom einfügt \cite{kadianakis2021tor}.

Der Netzwerksimulator Shadow ermöglicht es, ein vollständiges, deterministisches \gls{tor}-Netzwerk auf einer einzelnen Maschine zu betreiben \cite{jansen2011shadow}.
In dieser Umgebung lassen sich Padding-Konfigurationen unter sonst identischen Bedingungen vergleichen.

Die Wirksamkeit dieses produktiv ausgelieferten Schutzes gegen moderne \gls{df}-Angriffe ist bisher nur eingeschränkt untersucht.
Bestehende Studien stützen sich überwiegend auf Aufzeichnungen aus dem Live-Netzwerk oder auf vereinfachte Cell-Trace-Modelle, wodurch eine kontrollierte Variation einzelner Konfigurationen erschwert wird.


Hardwareanforderung einer kleineren Simulation war jetzt nicht ganz klar ?


Im Verlauf der Literaturrecherche zu dieser Arbeit rückten zusätzlich stärker eingreifende Schutzverfahren wie Tamaraw in den Fokus, die den Verkehrsstrom in feste Zeit- und Grössenraster zwingen \cite{cai2014systematic}.

\section{Fragestellung}

Die ursprüngliche Aufgabenstellung (Anhang~\ref{app:aufgabenstellung}) zielt auf den Aufbau einer Shadow-Simulation und die Auswertung des Tor-internen Circuit Padding.
Tamaraw wurde erst während der Erarbeitung dieser Arbeit als naheliegender Vergleichspunkt identifiziert und in Absprache mit dem Auftraggeber als ergänzender, exploratorischer Teil aufgenommen.
Daraus ergeben sich die folgenden Forschungsfragen.

\begin{enumerate}
    \item Wie verändert sich die Erkennungsrate eines \gls{df}-Angriffs unter den drei in \gls{tor} produktiv konfigurierbaren Modi des Circuit Padding (deaktiviert, voll aktiviert, reduziert), gemessen in einer Shadow-Simulation auf der \gls{hslu}-Infrastruktur?
    \item In welchem Verhältnis stehen die in der Simulation gemessenen Effekte zu den in der Literatur dokumentierten Limitationen simulierter und laborbasierter \gls{wf}-Studien?
    \item Welche zusätzliche Schutzwirkung lässt sich in derselben Simulationsumgebung durch das stärker eingreifende Tamaraw-Verfahren erzielen, und mit welchem Bandbreitenaufschlag ist sie verbunden?
\end{enumerate}

\section{Vision und erwartete Resultate}

Ziel der Arbeit ist es, den Einfluss von Circuit Padding und ergänzend von Tamaraw auf \gls{df}-Angriffe unter kontrollierten Laborbedingungen empirisch zu quantifizieren und die Ergebnisse in den Stand der Forschung einzuordnen.
Im Einzelnen werden folgende Resultate angestrebt.

\begin{itemize}
    \item Eine reproduzierbare Shadow-Simulationsumgebung in einer Linux-VM der \gls{hslu}-Infrastruktur, die ein verkleinertes \gls{tor}-Netzwerk inklusive Hintergrundlast abbildet.
    \item Ein in dieser Umgebung erhobenes \gls{wf}-Datenset aus \glspl{pcap} der Eingangsknoten, das für die drei Padding-Konfigurationen vergleichbar ist und sowohl Closed-World- als auch Open-World-Szenarien abdeckt.
    \item Ein mit der Bibliothek WFlib durchgeführter \gls{df}-Angriff, dessen Erkennungsrate über Padding-Modi, Klassenzahl und Szenario hinweg ausgewertet wird.
    \item Eine ergänzende, exploratorische Messreihe zum Tamaraw-Verfahren in derselben Umgebung, die Schutzwirkung und Bandbreitenaufschlag gegenüberstellt.
    \item Eine Einordnung der gemessenen Effekte vor dem Hintergrund dokumentierter Limitationen simulierter \gls{wf}-Studien.
\end{itemize}

\section{Aufgabenstellung}

Die vollständige, vom Auftraggeber unterzeichnete Aufgabenstellung findet sich in Anhang~\ref{app:aufgabenstellung}.
Sie umfasst den Aufbau der Shadow-Simulation und die Auswertung des Circuit Padding.
Die exploratorische Erweiterung um Tamaraw wurde während der Arbeit in Absprache mit dem Auftraggeber ergänzt.
